% apore-lite condensed study notes.
% Generated from this chapter's wiki/. Safe to edit by hand.
% Keep both APORE BODY markers exactly as they are:
% shared/scripts/build_notes.py extracts the text between them when
% building the combined domain PDF.
\documentclass[11pt]{article}
\usepackage{apore-notes}

\notesmeta{Fundamentals of Engineering Materials}{2. Atomic Structure and Interatomic Bonding}{2026-09-19}

\begin{document}
\makenotestitle
% >>> APORE BODY START
\chapterhead{Atomic Structure and Interatomic Bonding}

\begin{keydef}{Source keys}
\raggedright
\begin{points}
\item \textbf{Pre-lecture}: \texttt{before lecture-Chapter 2 - Atomic Structure
and Interatomic Bonding.pdf}
\item \textbf{Full deck}: \texttt{F26-D2L-Chapter 2- Introduction to
Engineering Materials-Part 1 and part 2 (full).pdf}
\item \textbf{Week 2 problems}: \texttt{EXAMPLE PROBLEMS 1 - WEEK 2.docx}
\end{points}
\end{keydef}

\topic{Atomic Structure}

\begin{keydef}{Particles}
Electrons ($9.11\times10^{-31}$ kg), protons and neutrons (each
$\approx 1.67\times10^{-27}$ kg). \src{Pre-lecture, slide 3}
\end{keydef}

\begin{keydef}{Atomic number}
The number of protons in the nucleus, equal to the number of electrons in a
neutral species. \src{Pre-lecture, slide 3}
\end{keydef}

\begin{keydef}{Quantum numbers}
\begin{points}
\item $n$, principal (shell): K, L, M, N, O for $1, 2, 3, 4, \dots$
\item $\ell$, azimuthal (subshell): s, p, d, f for $0, 1, 2, 3, \dots, n-1$
\item $m_\ell$, magnetic (number of orbitals): $1, 3, 5, 7$, ranging $-\ell$ to
$+\ell$
\item $m_s$, spin: $\tfrac{1}{2}$, $-\tfrac{1}{2}$
\end{points}
\src{Pre-lecture, slide 5}
\end{keydef}

\begin{keydef}{Maximum electrons per subshell}
$\mathrm{s}=2$, $\mathrm{p}=6$, $\mathrm{d}=10$, $\mathrm{f}=14$.
\src{Full deck, slide 22}
\end{keydef}

\begin{keydef}{Valence electrons}
Electrons in the outer unfilled shells. They are available for bonding and tend
to determine chemical properties. Filled shells are more stable and require more
energy to gain or lose electrons. For most elements the configuration is not
stable, because the valence shell is usually not completely filled.
\src{Pre-lecture, slide 8; Full deck, slide 23}
\end{keydef}

\begin{method}{Writing an electron configuration (Aufbau)}
\begin{steps}
\item Fill orbitals bottom-up, so electrons occupy the lowest available energy
states.
\item Respect the per-subshell maxima: s 2, p 6, d 10, f 14.
\item Identify the valence electrons as those in the outermost unfilled shell.
\end{steps}
\src{Full deck, slides 20--22}
\end{method}

\begin{worked}{Nitrogen}
Bottom-up filling gives $1s^2\,2s^2\,2p^3$, also written
$[\mathrm{He}]\,2s^2\,2p^3$. \src{Full deck, slide 22}
\end{worked}

\begin{keydef}{Reference configurations}
\begin{points}
\item Stable: He $1s^2$; Ne $1s^22s^22p^6$; Ar $1s^22s^22p^63s^23p^6$;
Kr $1s^22s^22p^63s^23p^63d^{10}4s^24p^6$
\item Others: H $1s^1$; Li $1s^22s^1$; Be $1s^22s^2$; B $1s^22s^22p^1$;
Na $[\mathrm{Ne}]3s^1$; Mg $[\mathrm{Ne}]3s^2$; Al $[\mathrm{Ne}]3s^23p^1$
\item C (6): $1s^2\,2s^22p^2$, valence $2s^22p^2$
\item Fe (26): $1s^2\,2s^22p^2\,3s^23p^6\,3d^64s^2$, valence $3d^64s^2$
\end{points}
\src{Pre-lecture, slides 7--9; Full deck, slide 25}
\end{keydef}

\begin{keydef}{Wave and particle character}
Electrons have both. Particle-like: an electron exists in a point orbiting a
nucleus. Wave-like: position in an orbital is a probability distribution, and
each orbital has a discrete energy level set by quantum numbers.
\src{Full deck, slide 18}
\end{keydef}

\topic{Atomic Mass, Atomic Weight, and the Mole}

\begin{keydef}{Approximate particle masses}
Proton $\approx 1$ amu, neutron $\approx 1$ amu, electron $\approx 0.00055$ amu.
$1\ \text{amu} \approx 1.67\times10^{-24}$ g $= 1.67\times10^{-27}$ kg.
\src{Full deck, slide 4}
\end{keydef}

\begin{keydef}{The amu}
$1$ amu $=\tfrac{1}{12}$ the mass of one carbon-12 atom, so one carbon-12 atom
$=12$ amu. Also written u. Grams are far too large for one atom, since a carbon-12
atom weighs only $\approx 1.99\times10^{-23}$ g. \src{Full deck, slide 9}
\end{keydef}

\begin{formula}{The amu--gram bridge}
\[
1\ \text{amu} = \frac{1\ \text{g}}{6.022\times10^{23}}
\approx 1.66\times10^{-24}\ \text{g}
\qquad
12\ \tfrac{\text{amu}}{\text{atom}} \;\longleftrightarrow\; 12\ \tfrac{\text{g}}{\text{mol}}
\]
\src{Full deck, slide 8}
\end{formula}

\begin{keydef}{The mole}
A counting unit, like a dozen: $1$ mol $=6.022\times10^{23}$ particles
(Avogadro's number), and $12$ g of carbon-12 $=1$ mol. A mole is a \emph{number}
unit like kilo or mega, so it carries no physical unit of its own.
\src{Full deck, slide 5; Week 2 problems, Problem 4}
\end{keydef}

\begin{keydef}{Isotopes}
Atoms of the same element with the same number of protons but a different number
of neutrons. Mass number $=$ protons $+$ neutrons; atomic number $=$ protons.
Carbon: $^{12}$C (6p, 6n, 98.93\%), $^{13}$C (6p, 7n, 1.07\%), $^{14}$C (6p, 8n,
trace). \src{Full deck, slide 12}
\end{keydef}

\begin{formula}{Atomic weight}
\[
\bar{A} = \sum (\text{isotope mass})(\text{fractional abundance})
\]
\src{Full deck, slide 14}
\end{formula}

\begin{worked}{Natural uranium}
$Z=92$, isotopes $99.274\%\ ^{238}$U, $0.720\%\ ^{235}$U, $0.006\%\ ^{234}$U.
Using mass numbers as approximate masses,
\begin{align*}
\bar{A}_{\mathrm{U}} &= (238)(0.99274) + (235)(0.00720) + (234)(0.00006)\\
&= 236.27212 + 1.69200 + 0.01404 = 237.97816\ \text{amu.}
\end{align*}
The published standard value is $238.02891$ amu; the gap comes from using
integer mass numbers instead of the actual isotope masses ($238.051$, $235.044$,
$234.041$). \src{Full deck, slide 14; Week 2 problems, Problem 1}
\end{worked}

\begin{formula}{Weight versus mass}
\[
W = m\,g
\]
Mass is the amount of matter (kg) and does not change; weight is the
gravitational force (N) and does. A 60 kg person: on Earth
$W=(60)(9.81)=588.6 \approx 589$ N; on the Moon $W=(60)(1.62)=97.2\approx 97$ N.
\src{Full deck, slide 6}
\end{formula}

\begin{pitfall}
Mass and weight are not the same quantity. Mass (kg) is the amount of matter
and stays the same, weight (N) is the gravitational force and changes with $g$.
\src{Full deck, slide 6}

Atomic mass and atomic weight are not interchangeable. Atomic mass refers to an
individual atom, atomic weight is the weighted average over naturally occurring
isotopes. \src{Full deck, slide 13}
\end{pitfall}

\topic{Periodic Table and Electronegativity}

\begin{keydef}{Electropositive and electronegative}
Electropositive elements readily give up electrons to become positive ions.
Group IA gives up 1, Group IIA gives up 2, the group left of the transition
metals gives up 3. Electronegative elements readily acquire electrons to become
negative ions, accepting 1 or 2 depending on group. Inert gases neither give up
nor accept. \src{Pre-lecture, slide 10}
\end{keydef}

\begin{keydef}{Electronegativity}
Ranges $0.7$ to $4.0$; large values indicate a tendency to acquire electrons. It
decreases toward the lower-left of the periodic table (Fr $=0.7$, Cs $=0.7$) and
increases toward the upper-right (F $=4.0$, O $=3.5$).
\src{Pre-lecture, slide 11}
\end{keydef}

\begin{keydef}{Values given}
H 2.1, Li 1.0, Be 1.5, C 2.5, N 3.0, O 3.5, F 4.0, Na 0.9, Mg 1.2, Cl 3.0,
Fe 1.8, Cu 1.9. \src{Pre-lecture, slide 11}
\end{keydef}

\topic{Bonding Forces and Energies}

\begin{formula}{Net energy and net force}
\[
E_N = E_A + E_R = -\frac{A}{r} + \frac{B}{r^n}, \qquad n \approx 8
\]
\[
F_N = F_A + F_R, \qquad F = -\frac{dE}{dr}, \qquad E = \int -F\,dr
\]
\src{Pre-lecture, slide 13; Full deck, slides 31, 44}
\end{formula}

\begin{formula}{The Coulombic constant $A$}
\[
A = \frac{1}{4\pi\varepsilon_0}\,\bigl(|Z_1|e\bigr)\bigl(|Z_2|e\bigr)
\]
with $\varepsilon_0 = 8.85\times10^{-12}$ F/m the permittivity of vacuum,
$|Z_1|,|Z_2|$ the absolute valences, and $e = 1.602\times10^{-19}$ C. This
assumes a totally ionic bond; since most bonds are not 100\% ionic, $A$ is
normally taken from experimental data rather than computed.
\src{Full deck, slide 44}
\end{formula}

\begin{keydef}{Equilibrium separation $r_0$}
The separation at which $E_N$ is a minimum ($E_0$) and $F_N = 0$. It is the bond
length, the distance between the centers of two bonded atoms at lowest
potential energy. For $r<r_0$ repulsion dominates and pushes atoms apart; for
$r>r_0$ attraction dominates and pulls them together; as $r\to\infty$ both are
negligible and $V\approx 0$. The equilibrium meant here is \emph{mechanical},
not chemical, not thermal.
\src{Pre-lecture, slide 13; Full deck, slides 32--35; Week 2 problems, Problem 3}
\end{keydef}

\begin{formula}{Lennard-Jones potential}
\[
V(r) = 4\varepsilon\left[\left(\frac{\sigma}{r}\right)^{12}
- \left(\frac{\sigma}{r}\right)^{6}\right]
\]
The $r^{-12}$ term is the strong short-range repulsion; the negative $r^{-6}$
term is the attraction. $\sigma$ and $\varepsilon$ are model parameters
determined by experiment. \src{Full deck, slides 36--38; Week 2 problems, Problem 3}
\end{formula}

\begin{keydef}{$\sigma$ and $\varepsilon$}
$\sigma$ is the separation where the potential crosses zero, $V(\sigma)=0$.
$\varepsilon$ is the depth of the potential well and is the bond energy, meaning the
energy needed to separate two bonded atoms from equilibrium to infinity. A
deeper well means stronger bonding. At equilibrium $V(r_0)=-\varepsilon$.
\src{Full deck, slides 34, 37, 38}
\end{keydef}

\begin{method}{Finding $r_0$ and the bond energy from Lennard-Jones}
\begin{steps}
\item Rewrite $V(r) = 4\varepsilon[\sigma^{12}r^{-12} - \sigma^{6}r^{-6}]$ and
differentiate:
$\dfrac{dV}{dr} = 4\varepsilon\left[-\dfrac{12\sigma^{12}}{r^{13}}
+ \dfrac{6\sigma^{6}}{r^{7}}\right]$.
\item Set $dV/dr = 0$ at $r=r_0$. Since $\varepsilon>0$, divide by $4\varepsilon$:
$\dfrac{12\sigma^{12}}{r_0^{13}} = \dfrac{6\sigma^{6}}{r_0^{7}}$.
\item Multiply by $r_0^{13}$: $12\sigma^{12} = 6\sigma^{6}r_0^{6}$, so
$r_0^{6} = 2\sigma^{6}$ and $r_0 = \sqrt[6]{2}\,\sigma \approx 1.122\,\sigma$.
\item Substitute back. With $(\sigma/r_0)^6 = \tfrac{1}{2}$ and
$(\sigma/r_0)^{12} = \tfrac{1}{4}$:
$E(r_0) = 4\varepsilon\left[\tfrac{1}{4}-\tfrac{1}{2}\right] = -\varepsilon$.
\end{steps}
\src{Full deck, slides 40--42; Week 2 problems, Problem 3}
\end{method}

\begin{worked}{Force between $\mathrm{K}^{+}$ and $\mathrm{O}^{2-}$}
Centers separated by 1.5 nm. The attractive force is the negative derivative of
the attractive energy with respect to separation, using
$A=\frac{1}{4\pi\varepsilon_0}(|Z_1|e)(|Z_2|e)$. Valences $+1$ and $-2$ give
$Z_1=1$, $Z_2=2$, and a magnitude of $2.05\times10^{-10}$ N, reported without
sign, since the force is known to be attractive.
\src{Full deck, slide 45; Week 2 problems, Problem 5}
\end{worked}

\begin{formula}{Bonding energy and melting temperature}
The larger the bonding energy $E_0$ (deeper well), the higher the melting
temperature $T_m$.

\begin{tabular}{@{}llrr@{}}
\toprule
\textbf{Bond} & \textbf{Material} & \textbf{kJ/mol} & \textbf{$T_m$ (\textdegree C)} \\
\midrule
Ionic    & NaCl              & 640  & 801 \\
         & LiF               & 850  & 848 \\
         & MgO               & 1000 & 2800 \\
         & CaF$_2$           & 1548 & 1418 \\
\addlinespace
Covalent & Cl$_2$            & 121  & $-102$ \\
         & Si                & 450  & 1410 \\
         & InSb              & 523  & 942 \\
         & C (diamond)       & 713  & $>3550$ \\
         & SiC               & 1230 & 2830 \\
\addlinespace
Metallic & Hg                & 62   & $-39$ \\
         & Al                & 330  & 660 \\
         & Ag                & 285  & 962 \\
         & W                 & 850  & 3414 \\
\addlinespace
van der Waals & Ar           & 7.7  & $-189$ \\
         & Kr                & 11.7 & $-158$ \\
         & CH$_4$            & 18   & $-182$ \\
         & Cl$_2$            & 31   & $-101$ \\
\bottomrule
\end{tabular}
\src{Pre-lecture, slide 29}
\end{formula}

\begin{pitfall}
$\sigma \neq r_0$. $\sigma$ is where the potential crosses zero; $r_0$ is where
it reaches its minimum, with $r_0 = 2^{1/6}\sigma$.
\src{Full deck, slides 34, 37}

The well depth is $+\varepsilon$ while the energy at the minimum is
$-\varepsilon$. Bond energy is quoted positive because it equals the positive
work required to break the bond.
\src{Full deck, slide 42; Week 2 problems, Problem 3}

\textbf{Source conflict.} The pre-lecture deck writes the Lennard-Jones
prefactor as $3\varepsilon$; the full deck and the Week 2 problems both use
$4\varepsilon$. \src{Pre-lecture, slide 30; Full deck, slides 36--38; Week 2 problems, Problem 3}
\end{pitfall}

\topic{Interatomic Bonding Overview}

\begin{keydef}{Primary and secondary}
Primary: ionic, covalent, metallic. Secondary (van der Waals, including
hydrogen bonding): from atomic or molecular dipole interactions. Secondary
bonding exists between virtually all atoms or molecules but may be obscured when
a primary type is present; it is weaker, yet still influences physical
properties. \src{Pre-lecture, slide 14; Full deck, slide 62}
\end{keydef}

\begin{keydef}{Directional versus non-directional}
\begin{points}
\item \textbf{Directional}: covalent and some secondary bonding. Electrons or
shared electrons are localized between atoms; the bond exists only because of
that location or dipole moment.
\item \textbf{Non-directional}: ionic and metallic. Electrons cannot be
localized; bonding forces are the same in all directions around the atom.
\end{points}
\src{Pre-lecture, slide 24}
\end{keydef}

\begin{formula}{Relative strength of intramolecular forces}
\[
\text{metallic} > \text{ionic} > \text{polar covalent} > \text{nonpolar covalent}
\]
\src{Pre-lecture, slide 27}
\end{formula}

\begin{keydef}{Intramolecular versus intermolecular}
Intramolecular forces hold atoms together within a molecule; intermolecular
forces bond molecules together. Boiling and melting points depend on the type
and strength of the intermolecular forces. Stronger forces require more energy
to break, so ionic and polar covalent compounds have higher boiling and melting
points. \src{Pre-lecture, slides 26, 28}
\end{keydef}

\begin{keydef}{Chapter summary}
Chemical, electrical, thermal and optical properties are determined by
electronic configuration. Materials with large bonding energies usually have
high melting temperatures. The nature of the bond depends on the electron
structure of the constituent atoms. \src{Full deck, slide 62}
\end{keydef}

\topic{Ionic Bonding}

\begin{keydef}{Ionic bonding}
Occurs between positive and negative ions. Requires electron transfer and a
large difference in electronegativity. A metal atom donates electrons and
becomes a stable cation; a non-metal accepts them and becomes a stable anion.
The two are held together by Coulombic (electrostatic) attraction.
\src{Pre-lecture, slides 15, 16}
\end{keydef}

\begin{worked}{MgO}
Mg ($[\mathrm{Ne}]3s^2$) donates its $3s^2$ electrons to O ($1s^22s^22p^4$),
forming $\mathrm{Mg}^{2+}$ ($[\mathrm{Ne}]$) and $\mathrm{O}^{2-}$
($[\mathrm{Ne}]$). \src{Pre-lecture, slide 16}
\end{worked}

\begin{keydef}{Where it appears}
The predominant bonding in ceramics. Examples: NaCl, MgO, CaF$_2$, CsCl.
\src{Pre-lecture, slide 18}
\end{keydef}

\begin{method}{Drawing a Lewis structure}
\begin{steps}
\item Write the element symbol.
\item Determine the number of valence electrons from the element's group.
chlorine is in group 7, so 7 valence electrons.
\item Draw dots for electrons around the symbol, starting at the right and going
counter-clockwise, pairing as needed.
\end{steps}
\src{Pre-lecture, slide 34}
\end{method}

\topic{Covalent Bonding}

\begin{keydef}{Covalent bonding}
Occurs between atoms with similar electronegativities, which share electrons.
Bonds involve valence electrons, normally s and p orbitals.
\src{Pre-lecture, slide 18}
\end{keydef}

\begin{keydef}{sp$^3$ hybridization}
Carbon can form sp$^3$ hybrid orbitals: an electron is promoted from the $2s$
orbital into the $2p$ orbital, then hybridizes, producing four equivalent
$2\mathrm{sp}^3$ orbitals at $109.5^\circ$ to one another.
\src{Pre-lecture, slide 19}
\end{keydef}

\begin{worked}{CH$_4$}
Carbon has 4 valence electrons and needs 4 more; each hydrogen has 1 and needs 1
more. Electronegativities of C and H are similar, so electrons are shared in
sp$^3$ hybrid covalent bonds, with orbital overlap between carbon's sp$^3$
orbitals and each hydrogen's $1s$ orbital. \src{Pre-lecture, slide 20}
\end{worked}

\topic{Metallic Bonding}

\begin{keydef}{Metallic bonding}
Electrons are delocalized to form an ``electron cloud.'' Positively charged ion
cores sit in a lattice surrounded by a sea of valence electrons delocalized
across the whole structure rather than bound to individual atoms. It is the
strongest of the intramolecular forces.
\src{Pre-lecture, slides 21, 27}
\end{keydef}

\topic{Secondary Bonding (Van der Waals)}

\begin{keydef}{Secondary bonding}
Arises from attractive forces between dipoles.
\begin{points}
\item \textbf{Fluctuating dipoles}: from asymmetric electron clouds; liquid
H$_2$ bonds between adjacent H$_2$ molecules this way.
\item \textbf{Permanent dipoles}: a molecule with a $+$ end and a $-$ end bonds
to another such molecule; liquid HCl, and chains in a linear polymer.
\end{points}
\src{Pre-lecture, slide 22}
\end{keydef}

\begin{keydef}{Hydrogen bonding}
A special type of secondary bonding where hydrogen is one of the constituents:
HF, H$_2$O, NH$_3$. Its magnitude is generally higher than other secondary
bonds, and the melting and boiling temperatures of those compounds are
correspondingly high. \src{Pre-lecture, slide 23}
\end{keydef}

\topic{VSEPR Theory and Molecular Geometry}

\begin{keydef}{VSEPR}
Valence Shell Electron Pair Repulsion: valence electron pairs around an atom
repel each other, and greater repulsion means higher energy and less stability.
VSEPR predicts the geometry that minimizes repulsion.
\src{Pre-lecture, slide 32}
\end{keydef}

\begin{method}{Predicting molecular geometry}
\begin{steps}
\item Draw the Lewis structure.
\item Count electron groups (bonding pairs and lone pairs) around the central
atom.
\item Determine the electron geometry.
\item Determine the molecular geometry from the number of bonding pairs versus
lone pairs.
\end{steps}
\src{Pre-lecture, slide 33}
\end{method}

\begin{keydef}{Notation $\mathrm{AX}_m\mathrm{E}_n$}
$A$ = central atom, $X$ = bonding atom, $E$ = non-bonded electrons, $m$ = number
of bonding atoms, $n$ = number of non-bonded electron pairs.
\src{Pre-lecture, slide 35}
\end{keydef}

\begin{formula}{Shape by formula type}
\begin{tabular}{@{}lll@{}}
\toprule
\textbf{Type} & \textbf{Geometry} & \textbf{Bond angle} \\
\midrule
AX$_2$ & planar          & $180^\circ$ \\
AX$_3$ & trigonal planar & $120^\circ$ \\
AX$_4$ & tetrahedral     & $109.5^\circ$ \\
\bottomrule
\end{tabular}
\src{Pre-lecture, slide 36}
\end{formula}

\begin{keydef}{Predicting distortions}
Repulsion strength, greatest to weakest: lone pair--lone pair $>$ lone
pair--bonding pair $>$ bonding pair--bonding pair. Size, largest to smallest:
lone pair $>$ triple bond $>$ double bond $>$ single bond. Unpaired valence
electrons take slightly more space than bonded electrons, which may shift bond
angles; double and triple bonds repel more strongly than single bonds.
\src{Pre-lecture, slides 38, 39}
\end{keydef}
% <<< APORE BODY END
\end{document}
